\documentclass[letterpaper,11pt]{article}

\usepackage{latexsym}
\usepackage[empty]{fullpage}
\usepackage{titlesec}
\usepackage{marvosym}
\usepackage[usenames,dvipsnames]{color}
\usepackage{enumitem}
\usepackage[hidelinks]{hyperref}
\usepackage{fancyhdr}
\usepackage[english]{babel}
\usepackage{fontawesome5}
\usepackage[scale=0.90,lf]{FiraMono}
\usepackage{contour}
\usepackage[normalem]{ulem}
\usepackage{tgheros}
\usepackage[T1]{fontenc}
\usepackage{microtype}

\definecolor{light-grey}{gray}{0.83}
\definecolor{dark-grey}{gray}{0.3}
\definecolor{text-grey}{gray}{.08}
\definecolor{accent-blue}{rgb}{0.12,0.31,0.47}
\definecolor{accent-teal}{rgb}{0.06,0.45,0.43}

\renewcommand{\ULdepth}{1.8pt}
\contourlength{0.8pt}
\newcommand{\myuline}[1]{%
  \uline{\phantom{#1}}%
  \llap{\contour{white}{#1}}%
}

\renewcommand*\familydefault{\sfdefault}

\pagestyle{fancy}
\fancyhf{}
\fancyfoot{}
\renewcommand{\headrulewidth}{0pt}
\renewcommand{\footrulewidth}{0pt}

\addtolength{\oddsidemargin}{-0.5in}
\addtolength{\evensidemargin}{0in}
\addtolength{\textwidth}{1in}
\addtolength{\topmargin}{-.5in}
\addtolength{\textheight}{1.0in}

\urlstyle{same}
\raggedbottom
\setlength{\tabcolsep}{0in}
\setlength{\parindent}{0pt}
\color{text-grey}

\titleformat {\section}{
    \color{accent-blue}\bfseries \vspace{2pt} \raggedright \large
}{}{0em}{}[\color{light-grey} {\titlerule[2pt]} \vspace{-4pt}]

\begin{document}

\begin{center}
    \textbf{\Huge \color{accent-blue} Arina} \\ \vspace{5pt}
    \small MSc Biochemistry, Jamia Millia Islamia, New Delhi
    \\ \vspace{3pt}
    \small \faEnvelope \hspace{2pt} \texttt{arinazubair01@gmail.com} \hspace{1pt} $|$
    \hspace{1pt} \faPhone \hspace{2pt} \texttt{+91 7906615993} \hspace{1pt} $|$
    \hspace{1pt} \faLinkedin \hspace{2pt} \href{https://www.linkedin.com/in/arina-430a2124b/}{\myuline{\texttt{linkedin.com/in/arina-430a2124b}}}
    \\ \vspace{7pt}
\end{center}

\begin{center}
  {\Large\bfseries\color{accent-blue} Motivation letter}
\end{center}
\vspace{-6pt}
{\color{light-grey}\noindent\rule{\linewidth}{2pt}}

\vspace{10pt}
4 May 2026

\vspace{10pt}
Dr.\ Stefanie Fruhwürth\\
Principal Investigator, Microglia and Neuroinflammation Research Group\\
Institute of Neuroscience and Physiology, Sahlgrenska Academy\\
University of Gothenburg, Sweden

\vspace{8pt}
\textbf{\color{accent-teal}Re: Application for PhD Position - PAR 2026/292 $|$ Identification of mechanisms and patterns of microglial activation in neurological disorders}

\vspace{10pt}
Dear Dr.\ Fruhwürth,

\vspace{6pt}
What if \textbf{HSV-1}, a virus most people carry without ever thinking about it, is quietly setting off a \textbf{microglial} cascade that, in the right person and given enough time, ends in \textbf{neurodegeneration}? I keep coming back to this. \textbf{HSV-1} reaches the brain, it dials down \textbf{TREM2}, it disrupts the \textbf{cGAS-STING} antiviral axis in \textbf{microglia}, and from there the response that was supposed to protect the brain starts going sideways. \textbf{Microglial activation} that should be short and targeted becomes chronic. Chronic activation becomes damage. And there is now growing evidence that herpesvirus exposure can act as an environmental nudge toward dementia in people who already carry \textbf{AD} risk in their immune pathways. Given that so much of the late-onset \textbf{AD} genetic risk sits in \textbf{microglial immune genes}, it starts to feel less like coincidence and more like the same cells that struggle to contain \textbf{HSV-1} may be the very ones driving \textbf{neurodegeneration} years down the line. Where the protective response tips into a damaging one, and what exactly flips that switch at the molecular level, is the question I really want to dig into during a PhD, and it is why I am applying to your group.

\vspace{6pt}
My immunology training began at Rapture Biotech International, where I received hands-on instruction in \textbf{ELISA}, rocket immunoelectrophoresis, radial immunodiffusion, \textbf{Western blotting}, and immunoprecipitation workflows. I developed a genuine feel for antibody-antigen specificity and for how assay design choices, blocking conditions, antibody concentrations, signal-to-noise calibration, affect the reliability of a readout. This matters directly for the \textbf{biomarker discovery} arm of your project: quantifying secreted \textbf{microglial markers} such as \textbf{sTREM2}, sCSF1R, or cytokines like \textbf{CXCL10} and \textbf{IL-6} in \textbf{CSF and blood} samples from \textbf{HSE} patients and \textbf{AD} cohorts is fundamentally \textbf{ELISA} work, and troubleshooting low-abundance analytes in clinical matrices is exactly the kind of problem where solid immunoassay instincts are useful. My interest in immune signalling at that stage also produced a co-authored abstract published in \textit{The Journal of Immunology} (2023, Vol.\ 210) on the SPHK1-S1P-S1PR1 pathway in endothelial barrier function, an early foray into how lipid-immune signalling shapes tissue-level inflammatory responses, which connects to the broader question of how \textbf{microglial activation states} are regulated and resolved.

\vspace{6pt}
At CSIR-IGIB, I carried out multi-marker \textbf{immunofluorescence} on embryonic and postnatal mouse brain tissue across a developmental time course. The pipeline: timed-mating crosses, embryo harvest at E12-E16, OCT embedding, cryosectioning, immunostaining with Sox2, Tbr2, and NeuN antibodies, \textbf{confocal} image acquisition, and quantitative analysis in \textbf{Fiji/ImageJ}, is directly transferable to the cellular imaging work in your lab. Characterising \textbf{microglial activation states} in \textbf{iPSC}-derived microglia after \textbf{HSV-1} infection, or in neuron-microglia co-cultures challenged with amyloid oligomers, involves the same core logic: multi-channel immunostaining for activation markers (Iba-1, \textbf{TREM2}, CD68, or phospho-STING), \textbf{confocal} imaging at appropriate resolution, and systematic quantification of morphological or fluorescence-intensity changes across conditions. My experience optimising antibody panels and confocal acquisition settings on heterogeneous tissue means I can contribute to these experiments from day one rather than learning the basics. I was acknowledged for this work in a publication in \textit{Developmental Biology} (Elsevier, 2026), DOI: \href{https://doi.org/10.1016/j.ydbio.2026.02.012}{\myuline{10.1016/j.ydbio.2026.02.012}}.

\vspace{6pt}
My \textbf{Western blotting} experience, developed across both the Rapture Biotech internship and my Master's thesis at ICGEB, maps onto a critical part of the mechanistic signalling work in this project. Studying how \textbf{HSV-1} suppresses \textbf{TREM2} and how that feeds into reduced \textbf{STING phosphorylation}, or how the \textbf{IKK$\varepsilon$-TBK1-IRF3} axis operates in \textbf{iPSC}-derived microglia versus other cell types, requires reliable immunoblotting for low-abundance phosphoproteins under infection conditions. At ICGEB, my thesis work involved iterating expression conditions, lysis buffers, and column chromatography to resolve a difficult \textit{Plasmodium falciparum} protein, and the discipline of optimising for a target that behaves unexpectedly is exactly what immunoblotting for activated kinases and cleaved effectors demands. My BRIC-THSTI internship additionally gave me experience with cDNA synthesis and RNA handling, which translates to \textbf{qRT-PCR} quantification of \textbf{microglial activation markers} (ISG induction, \textbf{TREM2} transcript changes, cGAS and STING expression) after experimental virus exposure.

\vspace{6pt}
\textbf{PCR-based genotyping}, which I performed routinely at IGIB for mouse colony management, also connects to the genetic component of this project. Screening \textbf{iPSC} lines derived from patients with defined genetic variants, such as the \textit{IKBKE} variant described in the \textit{JCI Insight} (2023) paper, or validating \textbf{CRISPR}-edited lines requires exactly this kind of careful PCR-based QC at each step. It is technically straightforward, but it is the kind of work where being systematic and practiced matters. Similarly, \textbf{qPCR} for viral genome quantification after \textbf{HSV-1} infection of \textbf{iPSC}-derived microglia is a natural extension of the genotyping workflows I already use.

\vspace{6pt}
As a supporting capability rather than my primary focus, my computational work in \textbf{single-cell RNA-seq} and \textbf{gene regulatory network} inference at Chalmers is relevant to the transcriptomic dimension of this project. Profiling \textbf{microglial activation states} after \textbf{HSV-1} infection or A$\beta$ treatment using scRNA-seq, and interpreting the resulting gene expression signatures in the context of known innate immune pathways, is something I can contribute to analytically. The \textbf{AD} GRN project I have been working on, which is under revision at \textit{npj Aging}, specifically involved identifying regulatory drivers in \textbf{microglia} and other glial cell types across Braak stages, directly relevant background for interpreting transcriptomic data from your \textbf{iPSC microglia} models.

\vspace{6pt}
I hold an MSc in Biochemistry from Jamia Millia Islamia (CGPA 9.39/10, Rank 2 of 30, First Division with Distinction) and the Gold Medal for the highest CGPA in my BSc Biotechnology cohort (CGPA 9.84/10). I attended the 17th Annual Conference of the International Association of Neurorestoratology (October 2025). My IELTS score is 7.5. The one clear gap is iPSC-based disease modelling, which I have not yet worked with, and that honesty is part of why I am applying. The combination of your lab's established iPSC microglia and neuron co-culture systems, the clinical collaboration with Prof.\ Zetterberg on CSF biomarkers, and the antiviral mechanistic work with Prof.\ Paludan's group is exactly the training environment I need to bring together the immunology bench skills I already have and the iPSC competence I want to build. I keep coming back to that opening question, and I think your lab is the right place to start properly answering it. I would be glad to talk about how my background fits what you are looking for.

\vspace{10pt}
Yours sincerely,

\vspace{16pt}
\textbf{Arina}\\
\small MSc Biochemistry $|$ CGPA 9.39/10 $|$ Gold Medallist, BSc Biotechnology

\end{document}
